\documentclass[11pt]{article}

\usepackage[margin=1in]{geometry}
\usepackage{amsmath}
\usepackage{booktabs}
\usepackage{enumitem}
\usepackage{graphicx}
\usepackage[table]{xcolor}
\usepackage{tabularx}
\usepackage[hidelinks]{hyperref}
\usepackage{microtype}

\definecolor{headerblue}{HTML}{1F4E78}
\definecolor{lightblue}{HTML}{DCE6F1}
\definecolor{lightgray}{HTML}{F2F2F2}

\hypersetup{
  pdftitle={Progress Report: FullNote and EHR XGBoost Baselines},
  pdfauthor={CS 7650 Project Team}
}

\setlist[itemize]{leftmargin=1.4em, itemsep=0.25em, topsep=0.35em}
\setlist[enumerate]{leftmargin=1.6em, itemsep=0.25em, topsep=0.35em}
\renewcommand{\arraystretch}{1.15}

\title{\textbf{Progress Report: FullNote and EHR XGBoost Baselines}\\
\large 30-Day Post-Discharge Mortality Prediction}
\author{CS 7650 Project Team}
\date{November 10, 2026}

\begin{document}

\maketitle

\begin{abstract}
This report summarizes the completed FullNote + electronic health record (EHR)
baseline for predicting 30-day post-discharge mortality. We implemented a fresh,
reproducible pipeline that validates the cohort, constructs structured EHR
features, embeds complete discharge summaries with frozen BioClinicalBERT, and
trains three XGBoost experiments. An EHR-only model provides the note ablation.
The direct-fusion experiment concatenates the EHR features and 768-dimensional
note embedding. The third experiment compresses the note
embedding into a leakage-safe mortality risk score before fusion with the EHR
features. On the held-out test set, adding the full note produced a modest
incremental improvement: AUROC increased from
0.8633 to 0.8890 and AUPRC from 0.3412 to 0.3771 relative to EHR only.
Risk-score fusion increased recall from 0.3146 to 0.4860 and
F1 from 0.3819 to 0.4010, but reduced discrimination and produced substantially
worse probability calibration. The direct model is therefore the stronger
general-purpose baseline, while the risk-score model illustrates a useful
sensitivity-oriented operating point.
\end{abstract}

\section{Objective and Current Scope}

The prediction target is a binary outcome for each eligible hospital admission:

\begin{quote}
\emph{Did the patient die within 30 days after hospital discharge?}
\end{quote}

The current work evaluates whether discharge notes add predictive value beyond
structured EHR data. We use the currently accepted cohort without regenerating
the data, and the previously identified hospice-related cohort issue is deferred.
All three experiments share the same validated cohort, EHR transformation, and
train/validation/test split. The two note experiments additionally share the
same frozen note embeddings and identity-alignment checks.

\section{Dataset}

The cohort is derived from MIMIC-IV structured EHR data and MIMIC-IV-Note
discharge summaries. The unit of prediction is an admission. Splits are separated
at the patient level so that a patient cannot occur in more than one split.

\begin{table}[ht]
\centering
\caption{Accepted cohort summary.}
\label{tab:cohort}
\begin{tabular}{lrrr}
\toprule
\rowcolor{lightblue}
Split & Admissions & 30-day deaths & Prevalence \\
\midrule
Train      & 31,792 & 1,599 & 5.03\% \\
Validation &  6,908 &   328 & 4.75\% \\
Test       &  6,782 &   321 & 4.73\% \\
\midrule
Total      & 45,482 & 2,248 & 4.94\% \\
\bottomrule
\end{tabular}
\end{table}

The full cohort contains 36,913 unique patients. Each admission includes stable
identifiers, the binary label, split assignment, admission-level demographics,
hospitalization attributes, diagnoses, ten measurement concepts, and one full
discharge summary. The pipeline validates the exact schema, unique admission and
note identifiers, timestamps, nonempty notes, binary labels, diagnosis content,
and patient-level split separation before feature construction.

\section{Completed Pipeline}

\subsection{Structured EHR features}

All learned preprocessing is fitted only on the training split and then reused
unchanged for validation and test. The resulting sparse EHR matrix has 4,041
features:

\begin{itemize}
  \item 63 numeric and measurement features, including age, length of stay,
        diagnosis count, measurement summaries, missingness indicators, and
        observation timing where available;
  \item 28 one-hot categorical features for gender, admission type, admission
        location, and insurance; and
  \item 3,950 sparse multi-hot diagnosis features from a training-only ICD
        vocabulary with a minimum training frequency of 10.
\end{itemize}

Identifiers, labels, split names, raw note text, and raw absolute timestamps are
not used as predictors.

\subsection{FullNote embeddings}

Each discharge summary is encoded with frozen BioClinicalBERT. Long notes are
split into non-overlapping chunks of at most 510 content tokens plus special
tokens. Padding and special tokens are excluded from pooling. Token hidden states
are summed across all chunks, divided by the note's total content-token count,
and L2-normalized to produce one 768-dimensional vector per admission.

The cached embedding artifact contains 45,482 rows and 768 columns. Generating it
required 350,762 chunks and 167,344,801 content tokens. A companion index stores
the subject, admission, and note identifiers, split, and label. Fusion joins by
note identifier and verifies every other identity field rather than assuming
that independently generated artifacts have the same row order.

\subsection{EHR ablation and two fusion experiments}

\begin{enumerate}
  \item \textbf{EHR only:} pass the 4,041 structured features directly to
        XGBoost without loading any note-derived feature.
  \item \textbf{Direct fusion:} concatenate 4,041 EHR features with the
        768-dimensional note vector, producing 4,809 predictors for XGBoost.
  \item \textbf{Risk-score fusion:} use logistic regression to map the note
        vector to one mortality score, append this score to the EHR matrix, and
        train XGBoost with 4,042 predictors.
\end{enumerate}

The risk-score experiment requires two distinct fitting roles. Five fold-specific
logistic models generate out-of-fold scores for training admissions, preventing
each training label from influencing its own note-derived feature. A final
logistic model fitted on all training embeddings generates validation, test, and
future inference scores. The five folds use stratified group splitting by patient:
all 31,792 training admissions receive exactly one out-of-fold score, and every
fold has zero patient overlap between its fitting and holdout partitions.

\section{Model Selection and Evaluation}

XGBoost hyperparameters and classification thresholds are selected using the
validation split only. For each model, the threshold is the validation cutoff
that maximizes F1. The test set remains untouched until final evaluation. This
threshold search changes only the final binary decision; threshold-independent
metrics such as AUROC and AUPRC are computed from the continuous model scores.

The EHR-only model selected \texttt{max\_depth=5},
\texttt{min\_child\_weight=1}, and 98 boosting rounds. Its
validation-selected threshold was 0.727023. The direct model selected
\texttt{max\_depth=7},
\texttt{min\_child\_weight=5}, and 1,177 boosting rounds. Its validation-selected
threshold was 0.261359. The risk-score model selected
\texttt{max\_depth=5}, \texttt{min\_child\_weight=5}, and 88 rounds, with a
threshold of 0.717312. In the latter model, the note-risk score ranked first among
all 4,042 inputs by XGBoost gain.

\section{Results}

\subsection{Note-only risk model}

Before EHR fusion, the logistic note-risk model showed consistent discrimination
across splits (Table~\ref{tab:note-risk}). The training result is out-of-fold and
therefore provides a leakage-safe estimate for the generated training feature.

\begin{table}[ht]
\centering
\caption{Note-only logistic risk-score performance.}
\label{tab:note-risk}
\begin{tabular}{lrrrr}
\toprule
\rowcolor{lightblue}
Split & AUROC & AUPRC & Brier score & Log loss \\
\midrule
Train (out-of-fold) & 0.7914 & 0.1611 & 0.0456 & 0.1783 \\
Validation          & 0.7981 & 0.1642 & 0.0429 & 0.1696 \\
Test                & 0.7855 & 0.1591 & 0.0429 & 0.1701 \\
\bottomrule
\end{tabular}
\end{table}

\subsection{Final test comparison}

\begin{table}[ht]
\centering
\caption{Held-out test performance at validation-selected thresholds.}
\label{tab:test-results}
\begin{tabular}{lrrr}
\toprule
\rowcolor{lightblue}
Metric & EHR only & Direct fusion & Risk-score fusion \\
\midrule
AUROC               & 0.8633 & \textbf{0.8890} & 0.8740 \\
AUPRC               & 0.3412 & \textbf{0.3771} & 0.3531 \\
Accuracy            & 0.9404 & \textbf{0.9518} & 0.9313 \\
Precision           & 0.3739 & \textbf{0.4856} & 0.3414 \\
Recall / sensitivity& 0.3832 & 0.3146 & \textbf{0.4860} \\
Specificity         & 0.9681 & \textbf{0.9834} & 0.9534 \\
F1                  & 0.3785 & 0.3819 & \textbf{0.4010} \\
Brier score $\downarrow$ & 0.1159 & \textbf{0.0381} & 0.1170 \\
Log loss $\downarrow$    & 0.3768 & \textbf{0.1689} & 0.3699 \\
\bottomrule
\end{tabular}
\end{table}

The test prevalence is 4.73\%, so accuracy must be interpreted cautiously: a
classifier predicting survival for every admission would achieve 95.27\%
accuracy while detecting no deaths. AUPRC is more informative for this imbalanced
task. The direct model's AUPRC of 0.3771 is approximately eight times the no-skill
prevalence baseline. Relative to EHR only, the full note improves test AUROC by
0.0256 and AUPRC by 0.0359, demonstrating additional threshold-independent
predictive signal.

At its selected threshold, EHR only identified 123 of 321 deaths, with 206 false
positives and 198 false negatives. Direct fusion identified 101 deaths, with 107
false positives and 220 false negatives. Risk-score fusion identified 156 deaths,
with 301 false positives and 165 false negatives. Thus, risk-score fusion found 55
additional deaths but generated 194 additional false alerts.

For the direct model, using the conventional 0.5 threshold instead of 0.261359
would increase precision from 0.4856 to 0.5873 and accuracy from 0.9518 to 0.9559,
but recall would fall from 0.3146 to 0.2305 and F1 from 0.3819 to 0.3311. The lower
validation-selected threshold is therefore appropriate when balanced precision
and recall, rather than accuracy alone, is the stated objective.

\section{Interpretation}

The EHR-only ablation confirms that the discharge summary adds predictive signal.
Direct fusion improves ranking, precision, specificity, accuracy, and probability
error relative to EHR only. Its independently selected threshold has lower recall,
which reflects an operating-point trade-off rather than worse ranking.

Direct fusion is the stronger overall baseline. It preserves the full note
representation and performs better on ranking quality (AUROC and AUPRC),
precision, specificity, and probability accuracy (Brier score and log loss).
Risk-score fusion compresses the note into a single supervised feature and loses
some information, but it shifts the operating point toward sensitivity and
achieves higher test recall and F1.

\subsection{Why the note improvement is modest}

A patient-level paired bootstrap with 1,000 resamples estimated the direct-fusion
minus EHR-only test AUROC change as $+0.0256$ (95\% interval $+0.0094$ to
$+0.0409$) and the AUPRC change as $+0.0359$ (95\% interval $-0.0005$ to
$+0.0745$). The AUROC gain is therefore likely real, but the AUPRC gain is less
precise and the selected-threshold F1 change is only $+0.0034$. The appropriate
conclusion is a modest incremental discrimination benefit, not a large
improvement in binary classification.

The implementation audit found no row-alignment, label, feature-width, inference,
or metric-reproduction error. All 45,482 note embeddings are finite,
768-dimensional, and L2-normalized. Nevertheless, the representations are highly
similar: across 20,000 random note pairs, median cosine similarity is 0.9868 and
mean similarity is 0.9856. The current encoder averages every non-special token
across a mean of 3,679 content tokens per discharge summary. Common clinical
language and boilerplate can therefore dominate the vector while short,
prognostically important sections are diluted. In addition, BioClinicalBERT is
frozen rather than optimized for mortality, and XGBoost receives the dense
embedding dimensions directly. These are limitations of the baseline design,
not evidence of a coding failure. Section-aware pooling, a task-specific note
head, or leakage-safe late fusion are natural next experiments.

The risk-score model's high threshold and large Brier score and log loss indicate
that its final XGBoost outputs should not be interpreted as calibrated absolute
mortality probabilities. Its scores are still useful for ranking and binary
classification at the empirically selected cutoff. If calibrated probabilities
are required, calibration must be added and assessed on data not used to fit the
base model.

\section{Engineering and Reproducibility}

The implementation was created as a separate package and does not import the
earlier project code. It provides command-line entry points for data validation,
EHR feature construction, note embedding, model training, saved-model inference,
and result plotting. Generated artifacts record the data contract, feature names,
model settings, thresholds, metrics, and identity indices. The embedding stage is
checkpointed and resumable, and saved inference reproduces the training-time
feature order and alignment checks. The current automated test suite contains 13
passing tests covering the data contract, feature construction, embedding
behavior, fusion alignment, EHR-only saved inference, and modeling utilities.

\section{Limitations}

\begin{itemize}
  \item The current cohort is accepted as provided. The hospice-related cohort
        concern is documented but intentionally not corrected in this iteration.
  \item Thresholds maximize validation F1; this is an experimental criterion, not
        a clinically validated operating policy.
  \item The report presents point estimates. Patient-level confidence intervals
        and formal statistical comparisons have not yet been computed.
  \item The experiment uses a single internal test split and has not been
        externally validated.
  \item The risk-score fusion output is not calibrated as an absolute
        probability, likely in part because the final classifier emphasizes the
        minority class.
\end{itemize}

\section{Next Steps}

\begin{enumerate}
  \item Select an operating threshold based on the desired trade-off between
        missed deaths and false alerts, and report performance at several fixed
        sensitivity targets in addition to maximum F1.
  \item Add patient-level bootstrap confidence intervals for AUROC, AUPRC, F1,
        recall, precision, Brier score, and the difference between experiments.
  \item Calibrate the risk-score fusion outputs using a held-out calibration
        procedure and re-evaluate Brier score, log loss, and calibration curves.
  \item Run targeted ablations, including note representation choices and EHR
        feature groups, to identify which components drive performance.
  \item Evaluate key demographic and admission subgroups for performance gaps.
  \item Preserve the direct-fusion model as the principal baseline and present
        risk-score fusion as the higher-recall alternative.
\end{enumerate}

\section*{Current Status}

The shared preprocessing and embedding pipeline, EHR-only ablation, both note
fusion experiments, saved inference paths, automated tests, metric summaries, and plots
are complete. The principal remaining work is evaluation refinement and report
integration rather than baseline implementation.

\end{document}
